
\subsection*{Abstract}
Federated Knowledge Graph reasoning requires nodes to identify relevant information across decentralized peers without compromising data privacy or bandwidth. We present \textbf{Holographic Indexing}, a two-tier discovery protocol designed for "Blind Semantic Search." Our method combines \textbf{Bloom Filter} \cite{bloom1970} membership probing for exact entity matching with \textbf{Community Centroid Signatures} for semantic neighborhood approximation. This "Holographic" representation allows nodes to identify potential reasoning paths in remote peers with minimal data leakage. We formalize the construction of these signatures and the \textbf{Synaptic Bridge Attention} mechanism that enables cross-graph traversals. We further describe security enhancements in v2.52.0, including \textbf{HMAC-SHA256 Path Provenance} \cite{gentry2009} and \textbf{Federated Leases}, which ensure the integrity and reliability of federated reasoning in high-velocity, multi-tenant environments. As of v2.52.0, full federated beam execution is realized through \texttt{DistributedBeamTraversal} and a dedicated \texttt{/traverse} endpoint enabling cross-node branch delegation, while the \texttt{ExternalValidator} (Phase 52) extends the federation concept to open scientific literature --- validating graph hypotheses against PubMed, arXiv, and OpenAlex.

\subsection*{1. Introduction}
The expansion of decentralized Knowledge Graphs necessitates a protocol for inter-graph discovery that respects the data sovereignty of individual nodes. Traditional federated search methods often require a central index or the exchange of full node lists, both of which are unacceptable in privacy-sensitive domains (e.g., healthcare or defense). Holographic Indexing addresses this by exchanging compressed, non-reversible topological signatures.

\subsection*{2. Tier 1: The Entity Hologram}
Each node $\mathcal{G}_i$ generates a Bloom Filter $B_i$ of its entity set $\mathcal{E}_i$. This allows a peer $\mathcal{G}_j$ to verify the existence of a specific entity $e$ with a configurable false-positive rate $p$, while ensuring that the full set $\mathcal{E}_i$ cannot be enumerated via the filter.

\subsection*{3. Tier 2: The Community Hologram}
To support fuzzy, semantic discovery, we introduce the Community Centroid Signature. For every community $C_k \in \mathcal{G}_i$, the node computes:
\begin{itemize}
  \item \textbf{Centroid}: $\vec{c}_k = \text{mean}(\{\vec{e} : e \in C_k\})$
  \item \textbf{Radius}: $\rho_k = \max \|\vec{e} - \vec{c}_k\|$
\end{itemize}


The set of all centroids forms the "Semantic Hologram." A remote reasoning beam looking for concept $\vec{x}$ can compute the "Synaptic Bridge Score":
$$\text{score}(C_k) = \cos(\vec{x}, \vec{c}_k)$$
Peers exceeding a threshold $\sigma$ (default 0.75) are flagged as relevant reasoning destinations.

\subsection*{4. Enterprise Security (v2.52.0)}
To prevent adversarial path injection in federated networks, v2.52.0 implements \textbf{HMAC-SHA256 Path Provenance}. Every reasoning response from a remote adapter is cryptographically signed using a shared secret $\mathcal{K}$. This ensures that "Synaptic Bridge" paths are both semantically valid and structurally authentic.

\subsection*{5. Conclusion}
Holographic Indexing provides a mathematically robust and privacy-preserving framework for federated graph reasoning. By decoupling exact membership from semantic relevance, it enables secure, decentralized intelligence at scale. In CEREBRUM v2.52.0, full federated beam execution is operational via \texttt{DistributedBeamTraversal} and the \texttt{/traverse} delegation endpoint, while the \texttt{ExternalValidator} extends the federation concept beyond peer KG nodes to open scientific literature, demonstrating that the Holographic discovery protocol generalizes naturally to heterogeneous federated knowledge sources.

\medskip\noindent\rule{\linewidth}{0.4pt}\medskip

\section*{6. Recent Advances (v2.51.1 -> v2.52.0)}

Federated reasoning has been one of the most actively developed areas of the CEREBRUM framework since v2.51.1. The following describes advances directly relevant to this paper.

\textbf{DistributedBeamTraversal and /traverse Endpoint (Phase 32).} The conceptual federation described at v2.51.1 is now fully implemented. \texttt{DistributedBeamTraversal} orchestrates a multi-node beam search where individual branches can be delegated to remote CEREBRUM nodes via the \texttt{/traverse} REST endpoint. Each remote node runs its own local beam segment and returns ranked partial paths; the coordinator merges and re-ranks these using the global CSA scoring function. This eliminates the need for a central index while preserving the beam's global coherence.

\textbf{RemoteAdapter Architecture.} A \texttt{RemoteAdapter} implements the standard \texttt{GraphAdapter} interface but routes all graph queries to a remote \texttt{/traverse} endpoint rather than a local data structure. This allows \texttt{DistributedBeamTraversal} to treat remote nodes as transparent graph backends, with Holographic Indexing determining which remote nodes are probed for each query.

\textbf{FederatedAdapter with Procrustes Embedding Alignment.} When multiple remote nodes use independently trained embedding spaces, cross-node semantic similarity scores are not directly comparable. The \texttt{FederatedAdapter} applies the same Orthogonal Procrustes alignment (PAPER_008) used for cross-modal signal encoding to align remote embedding spaces to a shared root space before computing Synaptic Bridge Scores. This prevents high-cosine false positives caused by unaligned embedding geometries.

\textbf{ExternalValidator: Scientific Literature Federation (Phase 52).} Beyond peer KG nodes, \texttt{ExternalValidator} extends the federation concept to open scientific databases. Given a hypothesis edge proposed by the \texttt{HypothesisEngine} (PAPER_007), \texttt{ExternalValidator} queries PubMed, arXiv, and OpenAlex for corroborating literature. Validated hypotheses receive elevated confidence scores; contradicted ones are suppressed. This represents a qualitative extension of "federated knowledge discovery" from structured graphs to unstructured scientific text.

\textbf{Security Enhancements.} HMAC-SHA256 path provenance, introduced at v2.51.1, is now enforced on all \texttt{/traverse} responses. Federated leases include TTL-bounded tokens to prevent replay attacks in long-running distributed sessions.